

\documentclass[final]{elsarticle}

\usepackage[margin=2.50cm]{geometry}  % tweak margins
%\usepackage{draftcopy}                % uncomment to print a "DRAFT-COPY" in background
\usepackage{graphicx}                 % if you want to include PostScript figures
\usepackage[figuresright]{rotating}   % if you have landscape tables
\usepackage{multirow}                 % fancy tables
\usepackage[toc,page]{appendix}       % appendix
\usepackage{lineno}                   % for printing the line numbers
\usepackage{natbib}                   % for building bibliography with \cite{}
\usepackage{setspace}
\usepackage{times}
\usepackage{afterpage}
\usepackage{float}
\usepackage{subcaption}
\usepackage{longtable}
\usepackage{blkarray}                 % for correlation matrices
\usepackage{color}
\usepackage{rotate}                   % rotate text in tables
\usepackage{lscape}                   % allow page to be landscape for large tables
\usepackage{amsfonts}                 % extra symbols
\usepackage{amsmath}
%\usepackage{floatpag}                 % \thisfloatpagestyle{empty}, to suppress page numbers
\usepackage{lmodern}
\usepackage{dcolumn}% Align table columns on decimal point
\usepackage{caption}% Use captionof to be able to minipage figure's \& tables together, but still get correct captions
\usepackage{enumitem}
\usepackage[breaklinks=true,hyperfigures=true]{hyperref} % make contents and references links, but do not colour them

\setenumerate{listparindent=\parindent} % Use paragraph indentation within the enumerate
\usepackage{cleveref}
% use to fix date
% \def \today { February 1st, 2009 }
% put your own definitions here:
\newcommand{\ck}{$\checkmark$}
\newcommand{\ttbs}{\char'134}
%\newcommand{\AmS}{{\protect\the\textfont2
%  A\kern-.1667em\lower.5ex\hbox{M}\kern-.125emS}}
% add words to TeX's hyphenation exception list
%\hyphenation{author another created financial paper re-commend-ed Post-Script}

\usepackage{listings}   
\lstset{language=C++}

% Particle aliases
\newcommand{\nue}{\nu_{e}}
\newcommand{\numu}{\nu_{\mu}}
\newcommand{\nutau}{\nu_{\tau}}
\newcommand{\nueb}{\xbar{\nu}_{e}}
\newcommand{\numub}{\xbar{\nu}_{\mu}}
\newcommand{\nutaub}{\xbar{\nu}_{\tau}}
\newcommand{\nub}{\xbar{\nu}}
\newcommand{\nua}{\nu_{\alpha}}
\newcommand{\nuab}{\xbar{\nu_{\alpha}}}

\newcommand{\numuany}{\overset{(-)}{\nu_\mu}}
\newcommand{\nueany}{\overset{(-)}{\nu_e}}


% By default, LaTeX doesn't like to fill more than 0.7 of a text page with tables and graphics, 
% nor does it like too many figures per page. This behaviour can be changed:
\renewcommand\floatpagefraction{.9}
\renewcommand\topfraction{.9}
\renewcommand\bottomfraction{.9}
\renewcommand\textfraction{.1}   
\setcounter{totalnumber}{50}
\setcounter{topnumber}{50}
\setcounter{bottomnumber}{50}

%define new fonts, smaller than tiny (for showing large covariance matrices in a page)
\makeatletter
\ifcase \@ptsize \relax% 10pt
  \newcommand{\miniscule}{\@setfontsize\miniscule{4}{5}}% \tiny: 5/6
  \newcommand{\verytiny}{\@setfontsize\miniscule{3}{4}}% \tiny: 5/6
\or% 11pt
  \newcommand{\miniscule}{\@setfontsize\miniscule{5}{6}}% \tiny: 6/7
  \newcommand{\verytiny}{\@setfontsize\miniscule{4}{5}}% \tiny: 6/7
\or% 12pt
  \newcommand{\miniscule}{\@setfontsize\miniscule{5}{6}}% \tiny: 6/7
  \newcommand{\verytiny}{\@setfontsize\miniscule{4}{5}}% \tiny: 6/7
\fi
\makeatother




\linenumbers

\begin{document}



{\Large \textbf{Comments on TN266v1.4 by Megan}} \\


--Replies (http://www.t2k.org/docs/technotes/266/Comments/reply\_Megan\_v1p0.pdf), line 3-30
$\rightarrow$ Did you add this info to the note somewhere?  Can you add it?  (I couldn't find it..)  I think it's useful.

\textcolor{blue}{DONE. Appendix E is added. Also caption of tables 18 - 21 has been improved.}


-- Replies, line 45-48
$\rightarrow$ Again, can you add this to the note? I couldn't find it.

\textcolor{blue}{Added in appendix E.}

-- Replies, line 128-129
$\rightarrow$ Does starting at the boundary of the fit range really remove possible
local minima problems in the fit?  I would think that would actually make it
worse, since the fitter can only start stepping in 1 direction... anyway it
may be worth it to think about expanding the fitter range to $\pm 2 \pi$ in future fits.

\textcolor{blue}{Since all the nuisance parameters are marginalized, i.e. the fit parameters are not
floating in n-dimensional space, using the method shown in section 5.4
the global minimum is found almost 100\% of the times.
However, when the number of free parameters minimized is increased (more than 2), 
fit strategy should be revisited and optimized.}

-- Replies, line 135-137
$\rightarrow$ Although the IH is excluded at $>90\%$ CL, T2K doesn't actually
have the sensitivity to actually exclude the IH (as was discussed at length last week). 
Then I think it's important to include the IH-only limits in the note as well.

\textcolor{blue}{DONE.}

-- TN266v1.4, line 780 + 783 + fig 20 : description doesn't match figure caption (IH vs NH).

\textcolor{blue}{DONE. Caption of fig 20 was wrong. Best-fit is inverted hierarchy, not normal. }

-- TN266v1.4, line 875-877, 878-879 : I'm suddenly worried about fitter biases pointing
our data toward Asimov Set A.  Line 878-879 -- this may be true now (and I'm not suggesting
delaying the release of this analysis), but since these things don't agree within error bars for
all parameter values, may need to re-think this for later versions of this analysis.
(I guess it's worse if things far away from our best fit point don't make sense, right?)..

\textcolor{blue}{The test values of $\sin^2 \theta_{23}$ are quite extreme, indeed 
when performing the posterior predictive method the critical values are close enough.
Also the difference when $\sin^2 \theta_{23}$ is fixed is not very large, though not fully negligible.
I think it's safe for this analysis release. However we need to make more studies for when 
higher statistic is collected.}

-- Table 27 caption : Asimov set A has $\delta_{CP} = -\pi/2$, right?

\textcolor{blue}{Yes it's correct. So the fake data sets were built by marginalizing all the oscillation parameters
and fixing $\delta_{CP}=0$ and normal hierarchy.
Then the fake data sets were weighted using the Asimov data set A, assuming it is the data set I'm fitting.}


\end{document}


